\section{Measures and Scoring}
\label{sec:scoring}
\subsection{Task Performance}
\paragraph{Active completion time.} The primary outcome will be elapsed active-search time in seconds from object reveal to the correct right-trigger selection, excluding the union of recorded pause intervals overlapping that period. Section~\ref{ch:measurement} specifies the clock equation and handling of overlapping pauses. Lower values will indicate faster completion. Readiness, opening guidance, post-selection feedback, questionnaires and between-trial breaks will not contribute to this outcome. An unfinished trial will retain its observed exposure and outcome, but will not receive a successful completion time. The full interval from opening-cue onset to selection will be retained separately to describe the complete instructed task duration.

\paragraph{Target contact and selection errors.} Time to first target gaze will be the active-clock time of the first valid, recorded gaze-ray intersection with the target, and time to first controller contact will use the first tracked right-controller ray--target intersection. Last target-gaze contact time will denote the last observed valid target intersection during the trial, not a fixation duration or a spoken confirmation. These event times will remain distinct from trigger selection and from the larger gaze-vicinity region used by the guidance system. Incorrect selections will be counted as trigger presses that select non-target objects. If no valid target contact is observed, its latency will remain missing rather than being set to zero, and all contact times will retain their source and sampling resolution.

\subsection{Starting Geometry and Search Behavior}
\paragraph{Starting angle, height and distance.} Before the opening cue, the system will record the reference origin, readiness-cross position and target position. Starting theta ($\theta_0$) will be the signed angle about world up between the horizontal origin-to-cross and origin-to-target vectors; its magnitude $|\theta_0|$ will quantify horizontal turn demand, and its sign will preserve direction. Degenerate horizontal vectors will be marked unavailable. The primary analysis will adjust for $|\theta_0|$ rather than divide completion time by degrees. Nominal wall-angle assignment, exact pre-cue theta and post-reveal head-to-target error will remain separate variables. The exact pre-cue snapshot is not yet implemented in the aligned study build, and existing nominal-theta fields are not substitutes for it; the snapshot must be implemented and verified before the study release is frozen.

Signed target-height difference will be $\Delta h=y_{\mathrm{target}}-y_{\mathrm{cross}}$, in metres; positive values will indicate above-cross targets. Euclidean target distance will be the three-dimensional distance between target center and the first valid tracked head position after reveal, with that observation's timestamp retained. This post-reveal distance will not substitute for a missing pre-cue geometry snapshot. All positions used in a calculation will share a recorded coordinate frame.

\paragraph{Head movement and zone coverage.} Gross head translation will sum distances between adjacent valid head positions, in metres, and gross head rotation will sum shortest quaternion angular displacements between adjacent valid orientations, in degrees. Fixed-axis component totals will be retained separately and will not be mislabeled as Euler yaw, pitch and roll. Invalid intervals, pauses and excessive sample gaps will break accumulation rather than be interpolated across, and covered time and gap counts will accompany totals. These measures describe head movement, not independently tracked chair rotation.

Distinct walls visited, wall transitions and revisits will summarize search coverage using stable valid gaze observations and a versioned zone-assignment rule. Transitions across tracking gaps will not be inferred. More motion or more visits will not by themselves establish worse strategy, since geometry and available observations also affect these totals.

\paragraph{Instructions and guidance timing.} The delivered active-search assistance total will count playback starts for directional guidance, gaze/pointing acknowledgments, and selection reminders, including delivered repeats; these categories will also be reported separately, as specified in Section~\ref{ch:measurement}. Requested repeats, delivered repeats and interrupted utterances will be retained separately; model requests, prefetched audio and success speech will not be counted as delivered search instructions. Opening cues and success feedback will retain separate event labels. Timing analysis will distinguish evidence qualification, decision-request latency, model deferral, scheduling delay and audio onset. Gaze confirmations and controller confirmations will be evaluated separately against their corresponding gates; the controller gate includes both its hold and the required nearby valid gaze.

\subsection{Subjective Measures}
\paragraph{Workload.} Six NASA-TLX dimensions---mental, physical and temporal demand, performance, effort and frustration---will be rated separately for each voice block during the final survey, after both blocks. The study will use a documented 0--10 adaptation without pairwise weighting, in which performance runs from Perfect to Failure and the other dimensions run from Very low to Very high. Individual ratings and the unweighted mean of all six completed dimensions will be retained, with higher scores indicating greater workload or poorer perceived performance. The source instrument is \href{https://doi.org/10.1016/S0166-4115(08)62386-9}{Hart and Staveland's NASA-TLX}.

\paragraph{Motivation and activity experience.} The \href{https://selfdeterminationtheory.org/intrinsic-motivation-inventory/}{Intrinsic Motivation Inventory} will contribute interest/enjoyment (7 items), effort/importance (5), pressure/tension (5), and value/usefulness (7). Responses will range from 1 = Not at all true to 7 = Very true. Published reverse keys will be applied before averaging each complete subscale. Interest/enjoyment will represent intrinsic motivation; the four subscales will not be combined into an overall IMI score. Activity wording will refer to the completed search block, with task-specific adaptations documented.

\paragraph{Agent perception.} Twenty-one items from \href{https://doi.org/10.1145/3651288}{Guo et al. (2024), Appendix A, Table A1} will assess co-presence (Q1--6), perceived intelligence (Q13--18), appealing/eerie/familiar impressions (Q20--22), and perceived anthropomorphism (Q37--42). ``Virtual character'' will be adapted to ``voice agent.'' Responses will range from 1 = Very strongly disagree to 7 = Very strongly agree. The three six-item constructs will have separate complete-scale means. Q20--22 will be analyzed individually, without a combined eeriness score. These attributed qualities will not be interpreted as actual agent intelligence or sentience.

\paragraph{Voice experience and reliance.} Twelve custom 1--7 agreement items will assess voice similarity, responsiveness, helpfulness, distraction, familiarity, trust, comfort, uncanniness, reliance, naturalness, willingness to follow and motivation to act without guidance. Voice similarity will serve as a manipulation check. Voice liking and retrospective percentage of guidance use will each use 0--100 responses with separate unable-to-answer options. These 14 numeric items will be analyzed separately; they are not a validated composite scale. Liking during voice audition will remain separate from condition-specific liking collected at the end.

\paragraph{Symptoms and debrief.} The 16 \href{https://doi.org/10.1207/s15327108ijap0303_3}{Simulator Sickness Questionnaire} symptoms will be rated from 0 = None to 3 = Severe before headset exposure and after headset removal. Canonical total or subscale scores will require the published membership and weights. Open debrief questions will explore perceived performance, guidance experience, own-voice comfort, possible use cases and voice preference; broad experience questions will precede specific eeriness or inner-dialogue probes.

Each condition-specific set in the final survey will therefore contain 65 numeric responses (130 across both sets): six workload ratings, 24 IMI items, 21 Guo items and 14 custom ratings. Technical checks and optional text will be additional. Questionnaire observations will have participant-by-block scope even when copied into per-trial summary files; they will not be treated as 40 independent responses. Missing answers will not be replaced with zero.


Table~\ref{tab:battery} summarizes the condition-specific battery. Counts refer to numeric items in one block-referenced set, not independent scales. Exact wording and reverse-key positions are preserved in Appendices~\ref{app:instruments} and~\ref{app:data}.
\begin{table}[htbp]
\centering
\caption[Condition-specific questionnaire battery]{Proposed final questionnaire battery for each referenced voice block. These ratings are administered after both blocks; setup voice ratings and the baseline/final SSQ are separate.}
\label{tab:battery}
\small
\tagpdfsetup{table/header-rows={1}}
\begin{tabularx}{\linewidth}{@{}p{0.25\linewidth}p{0.09\linewidth}X@{}}
\toprule
Measure & Items & Scoring and interpretation \\\midrule
NASA-TLX adaptation & 6 & Six 0--10 dimensions; complete unweighted mean. Preserve explicit bounds and adverse performance direction. \\
IMI & 24 & Separate interest/enjoyment (7), effort/importance (5), pressure/tension (5), and value/usefulness (7) means. No overall total. \\
Agent perception & 21 & Three six-item means: co-presence, intelligence, anthropomorphism. Appealing, eerie, and familiar remain three single items. \\
Custom voice and assistance & 14 & Individual ratings of voice and guidance experience. No unvalidated combined score. \\\midrule
Total per block reference & 65 & 130 numeric ratings across two references, before technical checks and optional text. \\\bottomrule
\end{tabularx}
\end{table}

\subsection{Measurement Quality}
Tracking validity, observation age, data coverage, gaps and spatial-validation error will accompany gaze-derived measures. Ray intersections and vicinity hits will not be labeled fixations, and render-frame observations will not establish native eye-tracker sampling frequency. Pupil-based workload inference and fixation/saccade outcomes will require independent validation before inclusion. Trial performance and subjective summaries will remain linked to separate raw gaze, motion and event records through participant, session, block and trial identifiers. Scale metadata will distinguish the present 0--10 workload form from any earlier exports using another range.

